\documentclass[10pt,a4paper]{amsart}
\usepackage[margin=19mm]{geometry}
\usepackage{amsmath,amssymb,mathtools}
\usepackage[scr=rsfs,cal=euler]{mathalfa}
\usepackage{xfrac}
\usepackage[unicode,hidelinks,bookmarks=false]{hyperref}
\newcommand{\ie}{i.e.\ }
\newcommand{\cf}{cf.\ }
\newcommand{\resp}{respectively\ }
\newcommand{\Subsection}{\subsection}
\providecommand{\nobreakdash}{\nobreak\hbox{-}}
\setlength{\emergencystretch}{2em}
\multlinegap0pt
\usepackage{bm}
\usepackage{bbm}
\usepackage{dsfont}
\usepackage{xspace}
\usepackage{cancel}
\usepackage{mathtools}
\usepackage{amsthm}
\makeatletter
\@ifundefined{thm}{\newtheorem{thm}{Theorem}}{}
\makeatother
\newcommand{\curly}{\mathrel{\leadsto}}
\title[Quantum Adams operations in quasimap K-theory]{Quantum Adams operations in quasimap K-theory}
\author{Shaoyun Bai, Jae Hee Lee}
\date{Source: arXiv:2510.09335v1 (CC BY 4.0)}
\begin{document}
\maketitle

\noindent\emph{Source and licence.} Quantum Adams operations in quasimap K-theory. Version arXiv:2510.09335v1 (CC BY 4.0), distributed under the Creative Commons Attribution 4.0 International licence, as recorded in the arXiv licence metadata for this exact version and shown on the abstract page https://arxiv.org/abs/2510.09335v1. Full rights notice: licenses/arxiv-2510.09335-CC-BY-4.0.txt.\par\smallskip

\noindent\textit{Editorial scope.} Counted unit: the Thm whose source label is \texttt{None} in the article (source0.tex lines 901--907, source label \texttt{None}), delivered together with the article's own complete original proof of it (source0.tex lines 908--939, one \texttt{proof} environment of 32 lines). No mathematics is written, repaired or re-derived: statement and proof are the article's own text line for line. Required context retained uncounted: none. Every definition, display and label of those lines is retained unchanged. Omitted from this package and not redistributed: 1--900, 940--2053. Nothing inside a retained excerpt is omitted or altered: every retained body is delivered exactly as the article's lines are written, so no omission receipt of any kind accompanies this package. Citations: the retained excerpts carry no citation command at all, so no bibliography entry of the article is transcribed and no reference is redistributed. Reference adaptation: none was needed and none was made; every reference inside the delivered unit resolves inside it, so no unresolved reference or citation is produced. Printed slips kept literal: no printed word, symbol, hypothesis or number of the counted statement or of its proof was altered. Layout and class: the article's own class is \texttt{amsart} with packages including inputenc, amssymb, mdframed, bm, bbm, xcolor, stmaryrd, tikz-cd, comment, xy, microtype, graphicx, amsthm, etoolbox; the wrapper is the lane's pinned \texttt{amsart} editorial wrapper at 10pt on A4, which restates the theorem-like environments the retained text needs and adds the article's own macro definitions that the retained text needs (\texttt{\textbackslash curly}), with extra wrapper packages (bm, bbm, dsfont, xspace, cancel, mathtools). The delivered numbering is the unit's own. Assets: no retained span contains a figure environment, a graphics-inclusion command, a PDF-page inclusion, a table or a TikZ picture; no image file is copied into this package, the package declares no asset dependency and no image file is redistributed with it. Licence: version arXiv:2510.09335v1 of this article is distributed under the Creative Commons Attribution 4.0 International licence, as recorded in the arXiv licence metadata for this exact version and shown on the abstract page https://arxiv.org/abs/2510.09335v1. A full rights notice is delivered at licenses/arxiv-2510.09335-CC-BY-4.0.txt. Characters: every retained excerpt is pure ASCII, so the delivered engine, XeLaTeX with its default Latin Modern text font, reports no missing character.
\begin{thm}
    For any $\mathcal{F} \in K_{\mathbf{T}}(X)[\![ z^{\mathrm{eff}} ]\!]$ and cocharacter $\sigma: \mathbb{C}^{\times} \to \mathbf{T}$, we have
    \begin{equation}
        Q\Psi^k_{\mathcal{F}}(q^{\sigma} a) \circ S_{\sigma} = S_{\sigma} \circ Q\Psi^k_{\mathcal{F}}(a),
    \end{equation}
    where $a$ denotes the equivariant variables.
\end{thm}

\begin{proof}
    The proof, as above, uses degeneration for an operator defined through counts of ($\sigma$-)twisted quasimaps. Consider the moduli space $\mathsf{QM}_d^{\sigma}(X)_{\mathrm{rel} \ p_1, p_2, p'}$, which carries relative evaluation maps
    \begin{equation}
        \mathrm{ev}_{p_1} \times \mathrm{ev}_k \times \mathrm{ev}_{p_2} : \mathsf{QM}_d^{\sigma}(X)_{\mathrm{rel} \ p_1, p_2, p'} \to E_\sigma \times E_\sigma^{\times k} \times E_\sigma
    \end{equation}
    where $E_\sigma \to \mathbb{P}^1$ is the fibration over $\mathbb{P}^1$ with fibers isomorphic to the stack $\mathfrak{X}$. We fix the equivariant structure so that $\mathbb{C}^\times_q$ acts trivially on the fiber over $p_2$, and acts by $\sigma$ on the fiber over $p_1$. The evaluation map at $p_i$ takes value at the fiber of $E_\sigma \to \mathbb{P}^1$ over $p_i \in \mathbb{P}^1$. We consider the relative compactification, hence the evaluation maps in fact land in the stable locus $X$.  Indeed, $\mathrm{ev}_k$ (considered as a map to $E_\sigma^k$) is $\mu_k$-equivariant with respect to the loop rotation of quasimaps and the cyclic permutation on the targets.
    
    Then one can consider the degeneration of the operator
    \begin{equation}
        \sum_ d z^d (\mathrm{ev}_{p_1} \times \mathrm{ev}_{p_2} )_*  \left(  \hat{\mathcal{O}}_{\mathrm{vir} ,\mathsf{QM}_d^{\sigma} (X)} \cdot   \mathrm{ev}_k^* (\mathbf{G}^{-1}\mathcal{F})_{eq}^{\boxtimes k} \right) \mathbf{G}^{-1}
    \end{equation}
    in two different ways. First we may degenerate the source curve $C$ at $p_1$ to  \begin{equation}
        C \curly C_1 \cup C'
    \end{equation}
    so that $C_1$ contains $p_1$ and the marked points $p_0', \dots, p'_{k-1}$ and maps to the fiber of $E_\sigma \to \mathbb{P}^1$ over $p_1 \in \mathbb{P}^1$ denoted by $X_1$, while $C'$ carries two marked points at the poles, one identified with $p_2$. The $\mu_k$-action rotates the two components simultaneously.

    The degeneration formula yields the operator
    \begin{equation}
       \sum_{d_1} z^{d_1} (\mathrm{ev}_{p_1} \times \mathrm{ev}_{p'} )_*  \left(  \hat{\mathcal{O}}_{\mathrm{vir} ,\mathsf{QM}_d (X_1)} \cdot   \mathrm{ev}_k^* (\mathbf{G}^{-1}\mathcal{F})_{eq}^{\boxtimes k} \right) \mathbf{G}^{-1}
    \end{equation}
    for the operator in $C_1$, which is exactly the configuration defining the quantum cyclic $k$th power operation, with the only difference being that $\mu_k$ does not act trivially on $X_1$. Since the identity map $X_1 \cong X$ with $X$ equipped with the trivial $\mu_k$-action is $\mu_k \times \mathbf{T}$-equivariant with respect to the group homomorphism $\mu_k \times \mathbf{T} \to \mu_k \times \mathbf{T}$ given by $(\eta, \theta) \mapsto (\eta, \theta + \sigma(\eta))$ for the cocharacter $\sigma$, this operator can be identified with $Q\Psi^k_\mathcal{F} (q^\sigma a)$.  The operator corresponding to $C'$-component is $S_\sigma$.

    Alternatively we may degenerate the source curve $C$ at $p_2$ to
    \begin{equation}
        C \curly C' \cup C_2
    \end{equation}
    where now $C_2$ contains $p_2$ and the marked point $p_0' , \dots, p'_{k-1}$ and maps to the fiber $X_2$ over $p_2 \in \mathbb{P}^1$. In this case, the operator corresponding to $C_2$ is exactly $Q\Psi^k_\mathcal{F}$, since the $\mu_k$-action on $X_2$ is already trivial. The operator corresponding to $C'$-component is $S_\sigma$.

    By comparing the two degenerations of the same operator, we arrive at the desired result.
    
    
\end{proof}

\end{document}
